\documentclass[9pt,a4paper,twocolumn]{ctexart}

% --- Packages ---
\usepackage[top=1.2cm, bottom=1.2cm, left=1.2cm, right=1.2cm, columnsep=0.8cm]{geometry}
\usepackage{hyperref}
\usepackage{graphicx}
\usepackage{float}
\usepackage{longtable}
\usepackage{tabularx}
\usepackage{booktabs}
\usepackage{enumitem}
\usepackage{xcolor}
\usepackage{fancyhdr}
\usepackage{listings}
\usepackage{fontspec}
\usepackage{titlesec}
\usepackage{caption}

% --- Code block styling ---
\definecolor{codebg}{RGB}{245,245,245}
\definecolor{codeframe}{RGB}{220,220,220}
\definecolor{codekw}{RGB}{0,0,192}
\definecolor{codecomment}{RGB}{0,128,0}
\definecolor{codestring}{RGB}{163,21,21}
\definecolor{codenum}{RGB}{128,0,128}
\definecolor{codepunct}{RGB}{90,90,90}
\definecolor{badgegreen}{RGB}{34,139,34}
\definecolor{badgeblue}{RGB}{30,144,255}
\definecolor{badgeblack}{RGB}{50,50,50}

% --- Difficulty badges (rounded corners via tikz, pre-rendered in saveboxes) ---
\usepackage{tikz}
\newsavebox{\badgechubox}
\savebox{\badgechubox}{\tikz[baseline]{\node[fill=badgegreen,text=white,rounded corners=2.5pt,inner xsep=4pt,inner ysep=1.5pt,font=\sffamily\footnotesize\bfseries]{初级};}}
\newsavebox{\badgezhongbox}
\savebox{\badgezhongbox}{\tikz[baseline]{\node[fill=badgeblue,text=white,rounded corners=2.5pt,inner xsep=4pt,inner ysep=1.5pt,font=\sffamily\footnotesize\bfseries]{中级};}}
\newsavebox{\badgegaobox}
\savebox{\badgegaobox}{\tikz[baseline]{\node[fill=badgeblack,text=white,rounded corners=2.5pt,inner xsep=4pt,inner ysep=1.5pt,font=\sffamily\footnotesize\bfseries]{高级};}}
\newcommand{\lvchu}{\usebox{\badgechubox}}
\newcommand{\lvzhong}{\usebox{\badgezhongbox}}
\newcommand{\lvgao}{\usebox{\badgegaobox}}
\pdfstringdefDisableCommands{%
  \def\lvchu{[初级]}%
  \def\lvzhong{[中级]}%
  \def\lvgao{[高级]}%
}

\lstset{
  basicstyle=\ttfamily\scriptsize,
  keywordstyle=\color{codekw}\bfseries,
  commentstyle=\color{codecomment}\itshape,
  stringstyle=\color{codestring},
  backgroundcolor=\color{codebg},
  frame=single,
  rulecolor=\color{codeframe},
  breaklines=true,
  breakatwhitespace=false,
  breakindent=0pt,
  postbreak=\mbox{\textcolor{red}{$\hookrightarrow$}\space},
  numbers=left,
  numberstyle=\tiny\color{gray},
  columns=flexible,
  keepspaces=true,
  showstringspaces=false,
  tabsize=2,
  aboveskip=4pt,
  belowskip=4pt,
  extendedchars=true,
  literate={，}{{，}}1 {；}{{；}}1 {！}{{！}}1 {？}{{？}}1 {“}{{“}}1 {”}{{”}}1 {（}{{（}}1 {）}{{）}}1 {《}{{《}}1 {》}{{》}}1,
}

% --- Extra language definitions (no built-in listings support) ---
\lstdefinelanguage{json}{
  morekeywords={true,false,null},
  sensitive=true,
  morestring=[b]",
  comment=[l]{//},
  morecomment=[s]{/*}{*/},
}
\lstdefinelanguage{yaml}{
  morekeywords={true,false,null,yes,no,on,off},
  sensitive=false,
  morecomment=[l]{\#},
  morestring=[b]",
  morestring=[d]',
}
\lstdefinelanguage{http}{
  morekeywords={GET,POST,PUT,DELETE,PATCH,HEAD,OPTIONS,HTTP,Host,Accept,Authorization,Content-Type,Content-Length,User-Agent},
  sensitive=true,
  morecomment=[l]{\#},
  morestring=[b]",
}
\lstdefinelanguage{TypeScript}{
  morekeywords={abstract,any,as,async,await,break,case,catch,class,const,continue,debugger,declare,default,delete,do,else,enum,export,extends,false,finally,for,from,function,get,if,implements,import,in,instanceof,interface,keyof,let,module,namespace,never,new,null,number,of,package,private,protected,public,readonly,return,set,static,string,super,switch,symbol,this,throw,true,try,type,typeof,undefined,unknown,var,void,while,with,yield,Record,Array,Promise,AsyncIterable,ReadableStream,Date,Error,Object,JSON},
  sensitive=true,
  morecomment=[l]{//},
  morecomment=[s]{/*}{*/},
  morestring=[b]",
  morestring=[b]',
  morestring=[b]`{`},
}

% --- Compact spacing ---
\linespread{0.95}
\setlength{\parskip}{1pt plus 1pt minus 1pt}
\setlength{\parindent}{0pt}

% --- Table styling ---
\renewcommand{\arraystretch}{1.1}

% --- Heading styling (numbered) ---
\titleformat{\section}{\normalsize\bfseries}{\thesection\quad}{0em}{}
\titlespacing{\section}{0pt}{6pt}{2pt}
\titleformat{\subsection}{\small\bfseries}{\thesubsection\quad}{0em}{}
\titlespacing{\subsection}{0pt}{4pt}{1pt}
\titleformat{\subsubsection}{\small\bfseries}{\thesubsubsection\quad}{0em}{}
\titlespacing{\subsubsection}{0pt}{3pt}{1pt}

% --- Page style ---
\pagestyle{fancy}
\fancyhf{}
\fancyhead[L]{\small\emph{\leftmark}}
\fancyhead[R]{\small\thepage}
\renewcommand{\headrulewidth}{0.4pt}
\renewcommand{\sectionmark}[1]{}  % prevent sections from overwriting leftmark

% --- Hyperref setup ---
\hypersetup{
  colorlinks=true,
  linkcolor=blue,
  urlcolor=blue,
  citecolor=blue,
}

% --- Caption format ---
\DeclareCaptionFormat{myformat}{\small\bfseries #1#2#3}
\captionsetup{format=myformat}

\begin{document}

\title{一文搞懂 Harness Engineering：六层架构、上下文管理与一线团队实战}
\date{}
\maketitle
\markboth{TypeScript 版 Agent 知识}{ TypeScript 版 Agent 知识}

\section*{Harness 基本概念}


\subsection*{Harness 到底是什么？}


可以先用一个粗暴但好记的说法：Agent = Model + Harness。你不是模型，那你做的东西大概率就是 Harness。



这个说法有点绝对，但抓住了重点。Harness 指的是模型之外的整套系统：系统提示词、工具调用、文件系统、沙箱环境、编排逻辑、钩子中间件、反馈回路、约束机制。模型只提供推理和生成能力，Harness 把状态、工具、反馈、执行环境和安全边界串起来，Agent 才能真正开始干活。



LangChain 的 Vivek Trivedi 写过一篇《The Anatomy of an Agent Harness》，里面有个思路很值得记：先分清模型负责什么，再看剩下的系统该补什么。用这条线一切，很多 Agent 问题就不再是“模型行不行”，而是“系统有没有把模型需要的东西准备好”。



可以把模型想成 CPU，把 Harness 想成操作系统。CPU 再强，OS 如果天天崩，体验也不会好。你买了最新的 M5 芯片，但系统卡死、驱动乱飞，实际体验可能还不如旧芯片配一个稳定系统。



\begin{figure}[H]
\centering
\includegraphics[width=0.85\textwidth]{/home/zhang/demos/latex_demos/Agent知识/HTML/一文搞懂 Harness Engineering：六层架构、上下文管理与一线团队实战 _ JavaGuide_files/harness-agent-equals-model-harness-arch.jpg}
\caption*{Agent = Model + Harness}
\end{figure}



\subsection*{Harness 和 Prompt / Context Engineering 的关系}


Prompt Engineering、Context Engineering、Harness Engineering 不太适合放在同一层比较。它们更像一层套一层，处理的问题范围越来越大。



\begin{figure}[H]
\centering
\includegraphics[width=0.85\textwidth]{/home/zhang/demos/latex_demos/Agent知识/HTML/一文搞懂 Harness Engineering：六层架构、上下文管理与一线团队实战 _ JavaGuide_files/harness-engineering-layers-arch.jpg}
\caption*{Harness 和 Prompt/Context Engineering 的关系}
\end{figure}


{\footnotesize
\begin{tabular}{|p{\dimexpr 0.232\columnwidth-2\tabcolsep\relax}|p{\dimexpr 0.207\columnwidth-2\tabcolsep\relax}|p{\dimexpr 0.256\columnwidth-2\tabcolsep\relax}|p{\dimexpr 0.305\columnwidth-2\tabcolsep\relax}|}
\hline
\textbf{层级} & \textbf{解决的问题} & \textbf{关注点} & \textbf{典型工作} \\
\hline
Prompt Engineering & 怎么把指令说清楚 & 让模型理解意图，减少局部歧义 & 系统提示词设计、Few-shot 示例、思维链引导 \\
\hline
Context Engineering & 该给 Agent 看什么 & 在合适时机给模型提供正确且必要的信息 & 上下文管理、RAG、记忆注入、Token 优化 \\
\hline
Harness Engineering & 系统怎么持续执行、纠偏、观测和恢复 & 长链路任务中的持续正确、偏差修正、故障恢复 & 文件系统、沙箱、约束执行、反馈回路、观测 \\
\hline
\end{tabular}
}



简单任务里，Prompt 可能就够了。比如让模型改一句文案，提示词说清楚，效果通常不会差。需要外部知识时，Context 更重要，你得把资料、检索结果、历史状态放到合适位置。到了长链路、可执行、低容错的商业场景，Harness 才会变成主要矛盾，因为 Agent 需要的不只是“会回答”，还要能执行、验证、回滚、继续推进。



这也是一线团队会把大量精力放在 Harness 上的原因。不是他们不会写 Prompt，而是 Prompt 解决不了所有执行问题。



\subsection*{Harness 包含哪些组件？}


想知道 Harness 里应该放什么，可以反过来问：模型做不到什么？



大模型看起来很能干，但从系统角度看，它仍然主要是一个输入输出函数。输入一段上下文，输出一段文本或结构化调用。它不会天然记住历史，不会自己跑命令，不会知道代码是否真的通过测试，也不会自动区分哪些信息该保留、哪些该丢掉。


{\footnotesize
\begin{tabular}{|p{\dimexpr 0.435\columnwidth-2\tabcolsep\relax}|p{\dimexpr 0.435\columnwidth-2\tabcolsep\relax}|p{\dimexpr 0.130\columnwidth-2\tabcolsep\relax}|}
\hline
\textbf{模型做不到的事} & \textbf{Harness 怎么补} & \textbf{对应组件} \\
\hline
记住多轮对话历史 & 维护对话历史，每次请求时拼进上下文 & 记忆系统 \\
\hline
执行代码、跑命令 & 提供 Bash 和代码执行环境 & 通用执行环境 \\
\hline
获取实时信息，比如新库版本、API 变化 & 接入 Web Search、MCP 工具 & 外部知识获取 \\
\hline
操作文件和环境 & 抽象文件系统，引入 Git 版本控制 & 文件系统 \\
\hline
判断自己有没有做对 & 提供沙箱、测试工具、浏览器自动化 & 验证闭环 \\
\hline
长任务中保持连贯 & 做上下文压缩、记忆文件、进度追踪 & 上下文管理 \\
\hline
\end{tabular}
}



把这些“模型做不了，但你又希望 Agent 能做到”的部分补齐，就是 Harness 的组件清单。LangChain 也把它拆成了几块：文件系统负责持久化，Bash 执行负责通用工具，沙箱负责隔离风险，记忆机制负责跨会话积累，上下文压缩负责对抗长上下文带来的质量下降。



\section*{Harness 进阶}


\subsection*{一个成熟的 Harness 长什么样？}


前面是从“模型缺什么，系统补什么”的角度看 Harness。如果换成系统设计视角，一个成熟的 Harness 通常会有清晰的分层。



我之前在 YouTube 上看到过一个六层体系，比较适合拿来理解 Harness 的全貌：



\begin{figure}[H]
\centering
\includegraphics[width=0.85\textwidth]{/home/zhang/demos/latex_demos/Agent知识/HTML/一文搞懂 Harness Engineering：六层架构、上下文管理与一线团队实战 _ JavaGuide_files/harness-engineering-six-layer-architecture.pdf}
\caption*{Harness Engineering 六层架构}
\end{figure}


{\footnotesize
\begin{tabular}{|p{\dimexpr 0.034\columnwidth-2\tabcolsep\relax}|p{\dimexpr 0.155\columnwidth-2\tabcolsep\relax}|p{\dimexpr 0.310\columnwidth-2\tabcolsep\relax}|p{\dimexpr 0.500\columnwidth-2\tabcolsep\relax}|}
\hline
\textbf{层级} & \textbf{名称} & \textbf{解决什么问题} & \textbf{关键设计} \\
\hline
L1 & 信息边界层 & Agent 该知道什么、不该知道什么 & 定义角色与目标，裁剪无关信息，结构化组织任务状态 \\
\hline
L2 & 工具系统层 & Agent 怎么和外部世界交互 & 选择工具、控制调用时机、提炼工具结果并反馈 \\
\hline
L3 & 执行编排层 & 多步骤任务怎么串起来 & 让模型按“理解目标、判断信息、分析、生成、检查”的轨道推进 \\
\hline
L4 & 记忆与状态层 & 长任务中间结果怎么管理 & 独立管理当前任务状态、中间产物和长期记忆，避免状态混在一起 \\
\hline
L5 & 评估与观测层 & Agent 怎么知道自己做对了没有 & 建立独立于生成过程的验证机制 \\
\hline
L6 & 约束、校验与恢复层 & 出错了怎么办 & 预设规则拦截错误，失败时提供重试、回滚或降级 \\
\hline
\end{tabular}
}



可以把它想成给一个新员工搭工作环境。L1 是岗位说明，告诉他该关注什么；L2 是办公工具；L3 是标准操作流程；L4 是项目管理系统和笔记本；L5 是质检流程；L6 是红线规则和应急预案。



这六层从边界、工具、流程、状态、验证到恢复，组成一整套体系。后面看 OpenAI、Anthropic、Stripe 的做法，会发现它们虽然形式不同，但很多设计都能映射到这六层。



不过不要一上来就想把六层全部搭齐。更现实的做法是先做 L1 和 L6：先让 Agent 知道自己该干什么，再给它设置出错后的拦截和恢复机制。这两层投入不算最高，但通常最容易见效。中间几层可以随着项目复杂度慢慢补。



\subsection*{为什么瓶颈经常不在模型？}


第一次听到这个结论，很多人会觉得反直觉。模型不够聪明，那等更强的模型出来不就好了？但不少实验和实践都在指向另一个结论：模型当然重要，但在很多 Agent 场景里，真正卡住效果的是基础设施。



前面提到的 \href{http://can.ac/}{Can.ac} 实验就是一个典型例子。同一个模型，只换了工具调用格式，效果能差十倍。LangChain 的实践也类似，他们优化了 Agent 运行环境，包括文档组织方式、验证回路、追踪系统，在 Terminal Bench 2.0 上从全球第 30 名升到第 5 名，得分从 52.8\% 提升到 66.5\%。模型没有换，换的是 Harness。



很多团队遇到 Agent 表现不好，第一反应是换模型或继续调提示词。这个反应很正常，但不一定命中问题。如果工具接口设计得很难用，反馈回路缺失，错误信息也不给修复方向，模型再强也会被外部环境拖住。



LangChain 还提到过一个 model-harness 耦合现象。现在很多 Agent 产品，比如 Claude Code、Codex，模型和 Harness 是一起被调优出来的，这会带来一种过拟合：模型习惯了某套工具逻辑，换一个 Harness 后表现可能变差。他们在 Terminal Bench 2.0 排行榜里观察到，Opus 在 Claude Code 的 Harness 下得分，远低于它在其他 Harness 中的得分。



他们的结论是：the best harness for your task is not necessarily the one a model was post-trained with。为任务选择 Harness 时，不要默认模型自带的 Harness 就一定最合适。



\subsection*{为什么上下文喂越多，Agent 反而越蠢？}


Dex Horthy 观察到一个很有意思的现象：168K token 的上下文窗口，用到大约 40\% 的时候，Agent 输出质量就开始明显下降。



\begin{figure}[H]
\centering
\includegraphics[width=0.85\textwidth]{/home/zhang/demos/latex_demos/Agent知识/HTML/一文搞懂 Harness Engineering：六层架构、上下文管理与一线团队实战 _ JavaGuide_files/context-utilization-40-percent-threshold-phenomenon.pdf}
\caption*{上下文利用率的 40\% 阈值现象}
\end{figure}


{\footnotesize
\begin{tabular}{|p{\dimexpr 0.278\columnwidth-2\tabcolsep\relax}|p{\dimexpr 0.222\columnwidth-2\tabcolsep\relax}|p{\dimexpr 0.500\columnwidth-2\tabcolsep\relax}|}
\hline
\textbf{区间} & \textbf{占比} & \textbf{表现} \\
\hline
Smart Zone & 0 - \textasciitilde\{\}40\\% & 推理聚焦、工具调用准确、代码质量高 \\
\hline
Dumb Zone & 超过 \textasciitilde\{\}40\\% & 幻觉增多、兜圈子、格式混乱、代码变差 \\
\hline
\end{tabular}
}



Anthropic 也遇到过类似问题，他们称之为“上下文焦虑”。Sonnet 4.5 在上下文快填满时会变得犹豫，甚至倾向于提前收工，即使任务还没完成。只做压缩不够，他们后来直接采用 context resets：清空上下文窗口，但通过结构化交接文档保留关键状态。



这里的目标不是给 Agent 塞更多信息，而是让它尽量停留在干净、相关的上下文里。一线团队做“渐进式披露”和“分层管理”，底层原因就在这里。上下文越多不等于越聪明，很多时候只是噪声越来越多。



生产环境里最好监控上下文利用率。一个可操作的做法是把 40\% 当成告警线，超过后触发压缩、分段执行或任务交接。等 Agent 已经开始兜圈子，再处理就比较被动了。



\subsection*{从哪里开始搭 Harness？}


结合一线团队的实践，可以把行动项按优先级拆开。没必要一开始做成大系统，先把 P0 做好，通常就能明显改善 Agent 表现。



\subsubsection{P0：可以马上做}

{\footnotesize
\begin{tabular}{|p{\dimexpr 0.256\columnwidth-2\tabcolsep\relax}|p{\dimexpr 0.436\columnwidth-2\tabcolsep\relax}|p{\dimexpr 0.308\columnwidth-2\tabcolsep\relax}|}
\hline
\textbf{行动} & \textbf{为什么} & \textbf{参考实践} \\
\hline
创建 \texttt{AGENTS.md} 并持续维护 & Agent 每次启动自动加载，犯错后更新，形成反馈循环 & Hashimoto 每一行对应一个历史失败案例 \\
\hline
写自定义 Linter + 修复指令 & 错误消息直接告诉 Agent 怎么改 & OpenAI 的 Linter 报错自带修复方法 \\
\hline
把团队知识放进仓库 & Slack、Wiki、Docs 里的知识对 Agent 很难稳定可见 & OpenAI 把仓库作为事实来源 \\
\hline
\end{tabular}
}



这里有个坑：不要把 \texttt{AGENTS.md} 写成超级 System Prompt。很多团队一上来恨不得把所有规则都塞进去，结果上下文被撑爆，Agent 反而更容易跑偏。OpenAI 的做法更克制，\texttt{AGENTS.md} 只当目录用，大约 100 行，详细规则放到子文档里按需加载。



\subsubsection{P1：P0 稳了之后再补}

{\footnotesize
\begin{tabular}{|p{\dimexpr 0.153\columnwidth-2\tabcolsep\relax}|p{\dimexpr 0.316\columnwidth-2\tabcolsep\relax}|p{\dimexpr 0.531\columnwidth-2\tabcolsep\relax}|}
\hline
\textbf{行动} & \textbf{为什么} & \textbf{参考实践} \\
\hline
分层管理上下文 & 避免把所有信息塞进一个文件，按需披露 & OpenAI 把 \href{http://agents.md/}{AGENTS.md} 当目录用，约 100 行 \\
\hline
建立进度文件和功能列表 & 用 JSON 追踪功能状态，Agent 不太容易乱改结构化数据 & Anthropic 初始化 Agent + 编码 Agent 两阶段 \\
\hline
给 Agent 端到端验证能力 & 让 Agent 像用户一样验证功能 & Anthropic 使用 Playwright / Puppeteer MCP \\
\hline
控制上下文利用率 & 尽量不超过 40\\%，用增量执行降低污染 & Dex Horthy 的 Smart Zone / Dumb Zone \\
\hline
\end{tabular}
}



\subsubsection{P2：有余力再考虑}

{\footnotesize
\begin{tabular}{|p{\dimexpr 0.164\columnwidth-2\tabcolsep\relax}|p{\dimexpr 0.463\columnwidth-2\tabcolsep\relax}|p{\dimexpr 0.373\columnwidth-2\tabcolsep\relax}|}
\hline
\textbf{行动} & \textbf{为什么} & \textbf{参考实践} \\
\hline
Agent 专业化分工 & 每个 Agent 携带更少无关信息，留在 Smart Zone & Carlini 的去重、优化、文档 Agent \\
\hline
定期垃圾回收 & 清理速度要跟得上生成速度 & OpenAI 的后台清理 Agent \\
\hline
可观测性集成 & 把性能优化从感觉问题变成可测量的问题 & OpenAI 接入 Chrome DevTools \\
\hline
\end{tabular}
}



\subsection*{你的 Harness 到哪个阶段了？}


可以用下面这个表粗略判断一下。这里不需要追求一步到 Level 4，很多团队能从 Level 0 到 Level 1，收益就已经很明显。


{\footnotesize
\begin{tabular}{|p{\dimexpr 0.304\columnwidth-2\tabcolsep\relax}|p{\dimexpr 0.464\columnwidth-2\tabcolsep\relax}|p{\dimexpr 0.232\columnwidth-2\tabcolsep\relax}|}
\hline
\textbf{阶段} & \textbf{特征} & \textbf{工程师角色} \\
\hline
Level 0：无 Harness & 直接给 Agent Prompt，没有结构化约束 & 手动写代码，偶尔使用 AI \\
\hline
Level 1：基础约束 & \texttt{AGENTS.md}、基础 Linter、手动测试 & 主要写代码，AI 辅助 \\
\hline
Level 2：反馈回路 & CI/CD 集成、自动化测试、进度追踪 & 规划和审查为主 \\
\hline
Level 3：专业化 Agent & 多 Agent 分工、分层上下文、持久化记忆 & 设计环境和管理执行过程 \\
\hline
Level 4：自治循环 & 无人值守并行化、自动清理、自修复 & 架构设计和质量把关 \\
\hline
\end{tabular}
}



\section*{Harness 还没解决的问题}


讲完这些实践，也要把没解决的问题摆出来。现在公开案例不少，但真正让人信服的方法论还不多，尤其是落到已有项目时，很多问题仍然悬着。


{\footnotesize
\begin{tabular}{|p{\dimexpr 0.102\columnwidth-2\tabcolsep\relax}|p{\dimexpr 0.259\columnwidth-2\tabcolsep\relax}|p{\dimexpr 0.639\columnwidth-2\tabcolsep\relax}|}
\hline
\textbf{问题} & \textbf{现状} & \textbf{谁在关注} \\
\hline
棕地项目怎么改造 & 公开成功案例几乎都是绿地项目，缺少成熟方法论 & Böckeler 把它比作“在从没用过静态分析的代码库上跑静态分析”。她还提出 Ambient Affordances：环境本身的结构特性，比如类型系统、模块边界、框架抽象，会影响 Harness 能做到什么程度 \\
\hline
怎么验证 Agent 做对了事 & 大家更擅长限制它别做错，但验证功能正确性还很弱 & Böckeler 批评：用 AI 生成的测试来验证 AI 生成的代码，仍然像“用同一双眼睛检查自己的作业” \\
\hline
AI 生成代码的长期可维护性 & LLM 代码经常重新实现已有功能，长期效果还不好判断 & Greg Brockman 提出过这个问题，但目前没有清晰答案 \\
\hline
Harness 该做厚还是做薄 & Manus 五次重写越做越简单，OpenAI 五个月越做越复杂 & 场景决定。通用产品更追求最小化，特定产品可以高度定制。模型变强后，已有 Harness 也应该定期简化，Anthropic 已经做过类似验证 \\
\hline
单 Agent 还是多 Agent & Hashimoto 坚持单 Agent，Carlini 使用 16 个并行 Agent & 规模决定。小项目单 Agent 往往够用，大项目更容易走向专业化分工 \\
\hline
\end{tabular}
}



绿地项目和棕地项目是软件工程里的经典说法。绿地项目指从零开始的新项目，没有历史包袱，就像在空地上盖房子，想怎么设计都比较自由。棕地项目指在已有代码库上改造，里面有历史架构、技术债和遗留逻辑，就像在老旧城区翻新，很多管线不能随便动。



OpenAI、Anthropic、Stripe、Hashimoto 这些案例基本都是在新项目里从零搭 Harness。但现实里，大多数团队面对的是跑了多年的老代码库。一个有十年历史、没有明确架构约束、到处是技术债的项目，怎么引入 Harness？目前还没有公开的成熟方法论。



\section*{Harness 案例：这些团队是怎么做的}


下面几个案例放在一起看，会发现不同背景的团队踩坑很像。区别主要在于，有的团队先撞墙再补 Harness，有的团队从第一天就把约束和反馈回路放进架构里。



\subsection*{OpenAI：三个人，五个月，一百万行，零手写代码}


先看数据：


{\footnotesize
\begin{tabular}{|p{\dimexpr 0.304\columnwidth-2\tabcolsep\relax}|p{\dimexpr 0.696\columnwidth-2\tabcolsep\relax}|}
\hline
\textbf{指标} & \textbf{数值} \\
\hline
团队规模 & 3 名工程师，后扩至 7 人 \\
\hline
持续时间 & 5 个月，2025 年 8 月起 \\
\hline
代码规模 & 约 100 万行 \\
\hline
手写代码 & 0 行，设计约束 \\
\hline
合并 PR 数 & 约 1,500 个 \\
\hline
日均 PR/人 & 3.5 个 \\
\hline
效率提升 & 约 10 倍 \\
\hline
\end{tabular}
}



数字很夸张，但更值得看的是他们怎么做。



\subsubsection{给 Agent 一张地图，不要塞一本千页手册}


OpenAI 的 \texttt{AGENTS.md} 大约只有 100 行，作用更像目录，指向 \texttt{docs/} 目录下更深层的设计文档、架构图、执行计划和质量评级。这就是渐进式披露：先给最关键的信息，需要更多细节时再加载。



这和到一个新城市很像。你不需要一上来背完整本旅游指南，先给一张地图，再告诉你想了解某个景点时去翻哪一页，就够用了。



Agent Skills 也可以看成渐进式披露的一种实现。它保留少量元数据，比如名称和描述，详细规则和执行流程只在触发时再加载进上下文。这个思路和 OpenAI 把 \texttt{AGENTS.md} 当目录很接近，只是 Skills 把这个模式标准化了。相关阅读：《Agent Skills 详解：是什么？怎么用？和 Prompt、MCP 有什么区别？》。



\subsubsection{架构约束要靠工具执行}


OpenAI 给每个业务领域定义了固定分层：



\begin{lstlisting}
Types → Config → Repo → Service → Runtime → UI
\end{lstlisting}



依赖方向不能反过来。怎么保证？靠自定义 Linter 和结构测试。违反规则时，工具不只是报错，还会告诉 Agent 应该怎么改。Agent 在修错的过程中，也被反复训练成更符合团队规范的写法。



OpenAI 有句原话很直接：If it cannot be enforced mechanically, agents will deviate. 只写在文档里的约束不够，不能机械化执行，Agent 迟早会偏离。



\subsubsection{可观测性也要给 Agent 看}


他们把 Chrome DevTools Protocol 接进 Agent 运行时，Agent 可以自己抓 DOM 快照和截图。日志、指标、链路追踪也通过本地可观测性栈暴露给 Agent。



这样一来，“把启动时间降到 800ms 以下”就变成了一个 Agent 可以自己测量、自己验证的目标。



\subsubsection{熵不会自己消失}


AI 生成代码越多，低质量实现、重复逻辑、文档不一致也会跟着变多。一开始 OpenAI 团队每周五花 20\% 时间手动清理这些生成物。后来这件事被自动化了：后台 Agent 定期扫描文档不一致、架构违规和冗余代码，并自动提交清理 PR。



这个点很现实。生成速度上来了，如果清理速度跟不上，项目迟早会被自己的产物拖垮。



\subsubsection{Slack 里的知识，Agent 很难稳定用上}


写在 Slack 讨论或 Google Docs 里的知识，对 Agent 来说并不稳定。OpenAI 的做法是把团队知识作为版本控制制品放进仓库里，让仓库成为可追踪、可引用的事实来源。



这里也别误解成“照抄 OpenAI 就行”。OpenAI 自己也说了，这个结果不应该被假设为在缺少类似投入的情况下可以复现。它的每一项方法都要前期投入。真正适合普通团队先学的，是地图式文档、机械化约束和主动清理这些思路。



\subsection*{Anthropic：从上下文焦虑到三智能体架构}


Anthropic 在这个方向上有两个值得细看的实践。一个是 Carlini 用多 Agent 写 C 编译器，另一个是 Anthropic Labs 借鉴 GAN 思路做三智能体协作。



\begin{figure}[H]
\centering
\includegraphics[width=0.85\textwidth]{/home/zhang/demos/latex_demos/Agent知识/HTML/一文搞懂 Harness Engineering：六层架构、上下文管理与一线团队实战 _ JavaGuide_files/anthropic-three-agent-collaborative-architecture-inspired-by-gan.pdf}
\caption*{Anthropic 三智能体协同架构（受 GAN 启发）}
\end{figure}



\subsubsection{用 16 个 Agent 写 C 编译器}


Nicholas Carlini 用大约两周时间，跑了 16 个并行 Claude Opus 实例，大约 2000 个 Claude Code 会话，做出了一个 GCC torture test 通过率 99\% 的 C 编译器。


{\footnotesize
\begin{tabular}{|p{\dimexpr 0.225\columnwidth-2\tabcolsep\relax}|p{\dimexpr 0.775\columnwidth-2\tabcolsep\relax}|}
\hline
\textbf{指标} & \textbf{数值} \\
\hline
持续时间 & 约 2 周 \\
\hline
并行 Agent 数 & 16 个 Claude Opus 实例 \\
\hline
会话数 & 约 2,000 个 \\
\hline
产出 & 10 万行 Rust 代码 \\
\hline
GCC torture test & 99\\% 通过率 \\
\hline
可编译项目 & PostgreSQL、Redis、FFmpeg、CPython、Linux 6.9 Kernel 等 150+ \\
\hline
API 成本 & 约 2 万美元 \\
\hline
\end{tabular}
}



这个项目里的 Harness 细节比结果本身更值得看：


\begin{itemize}[itemsep=1pt, topsep=2pt]
  \item 日志不打到控制台，全部写进文件，并使用 grep 友好的单行格式，比如 \texttt{ERROR: [reason]}，主动减少上下文污染。
  \item 测试不全部跑。每个 Agent 只跑随机 1-10\% 的测试子集；对单个 Agent 来说，子采样是确定性的，同一次运行总是跑同样的子集；跨 VM 又是随机的，不同 Agent 覆盖不同部分。这样整体覆盖全部测试，单个 Agent 不会在测试上耗掉几个小时。
  \item Agent 角色逐渐专业化，包括核心编译器工作、去重、性能优化、代码质量和文档。LLM 经常重新实现已有功能，所以专门做去重也很有必要。
\end{itemize}


Carlini 后来说过一句话：”我必须不断提醒自己，我是在为 Claude 写这个测试框架，不是为自己写。”这句话点出了 Harness 的服务对象：首先是 Agent，不一定是人类工程师。



\subsubsection{Anthropic 为什么借鉴 GAN？}


Anthropic Labs 团队在 2026 年 3 月发布了一个受 GAN 思路启发的三智能体架构。原文说的是 Taking inspiration from GANs，意思是借鉴思路，并不是真正做对抗训练。



\begin{lstlisting}
Planner（规划者）→ Generator（执行者）⇄ Evaluator（评估者）
\end{lstlisting}



Planner 拿到 1-4 句话的产品描述，把它扩展成完整产品规格，并被要求“在范围上要大胆”。Generator 按功能一个个做 Sprint，每个 Sprint 有明确完成标准。Evaluator 用 Playwright MCP 实际点击运行中的应用，再按产品设计深度、功能性、视觉设计、代码质量等维度打分。



这个架构主要处理两个问题：


{\footnotesize
\begin{tabular}{|p{\dimexpr 0.103\columnwidth-2\tabcolsep\relax}|p{\dimexpr 0.397\columnwidth-2\tabcolsep\relax}|p{\dimexpr 0.500\columnwidth-2\tabcolsep\relax}|}
\hline
\textbf{问题} & \textbf{表现} & \textbf{解法} \\
\hline
上下文焦虑 & Sonnet 4.5 快到上下文上限时草草收尾 & context resets + 结构化交接，单靠压缩不够 \\
\hline
自我评价偏差 & Agent 自信地夸自己做得好，实际质量一般 & 生成和评估交给两个独立 Agent \\
\hline
\end{tabular}
}



打分标准也有意思。前端设计里，设计质量和原创性的权重被故意调得比功能性和代码质量更高，因为模型很容易做出“功能齐全但长相平庸”的东西。权重调整是在逼它往更难的方向走。



\subsubsection{遇到上下文焦虑，Anthropic 选择重启}


Anthropic 发现 Sonnet 4.5 在上下文快满时会变得犹豫，甚至提前收工。他们最后采用的方案叫 context resets。



流程很简单：当 Agent 上下文接近饱和时，先把当前任务状态、已完成工作、待办事项结构化提取出来；然后启动一个新的干净 Agent，把交接文档给它；新 Agent 从干净状态继续做。



这有点像程序遇到内存泄漏。你不一定非要手动释放每个内存块，也可以重启进程，再从检查点恢复状态。听起来粗暴，但长任务里，一个干净的新 Agent 往往比一个塞满历史信息的 Agent 表现更好。



这个思路和 Carlini 的编译器项目也很接近。他跑了 2000 个 Claude Code 会话，每个会话都相对独立，从干净状态开始。Anthropic 只是把“重启和恢复”做得更正式。



两种配置的成本对比如下：


{\footnotesize
\begin{tabular}{|p{\dimexpr 0.622\columnwidth-2\tabcolsep\relax}|p{\dimexpr 0.111\columnwidth-2\tabcolsep\relax}|p{\dimexpr 0.089\columnwidth-2\tabcolsep\relax}|p{\dimexpr 0.178\columnwidth-2\tabcolsep\relax}|}
\hline
\textbf{配置} & \textbf{耗时} & \textbf{花费} & \textbf{效果} \\
\hline
Solo Harness，单 Agent + 最少工具 & 20 分钟 & \$9 & 跑不起来的半成品 \\
\hline
Full Harness，三 Agent + 完整工具链 & 6 小时 & \$200 & 完整可用的应用 \\
\hline
\end{tabular}
}



更复杂的任务差距还会拉大。比如用 Full Harness 做一个浏览器里的音乐制作工作站 DAW，跑了将近 4 小时，花了 \$124.70，最后得到一个带编曲视图、混音台和播放控制的可用程序。



但他们还有一个重要发现：把模型从 Sonnet 4.5 换成 Opus 4.6 后，Sprint 机制可以完全移除，Evaluator 从每个 Sprint 检查变成最后只检查一次。Anthropic 的总结很准确：Every component in a harness encodes an assumption about what the model can't do on its own, and those assumptions are worth stress testing.



即，Harness 里的每个组件都在假设"模型自己做不到这个"。模型变强后，这些假设要重新测试。Anthropic 也提到，模型越强，Harness 的设计空间会移动，旧的保护机制可能会变成冗余，所以 Harness 也要定期简化。



\subsection*{Stripe：每周 1300+ 个 PR 的无人值守模式}


Stripe 的 Minions 系统是另一个极端：高度自动化、无人值守。开发者发一条 Slack 消息，Agent 就从写代码、跑 CI 到提 PR 全部完成，人只在最后审查。每周有超过 1300 个完全由 Minions 生产、没有人类手写代码的 PR 被合并。



\begin{figure}[H]
\centering
\includegraphics[width=0.85\textwidth]{/home/zhang/demos/latex_demos/Agent知识/HTML/一文搞懂 Harness Engineering：六层架构、上下文管理与一线团队实战 _ JavaGuide_files/stripe-hybrid-state-machine-orchestration-architecture.pdf}
\caption*{Stripe 混合状态机编排架构}
\end{figure}



这个数字第一次看到确实有点吓人。拆开看，它靠的是一套很成熟的工程环境，不是某个”超强 Agent”。


{\footnotesize
\begin{tabular}{|p{\dimexpr 0.156\columnwidth-2\tabcolsep\relax}|p{\dimexpr 0.052\columnwidth-2\tabcolsep\relax}|p{\dimexpr 0.792\columnwidth-2\tabcolsep\relax}|}
\hline
\textbf{组件} & \textbf{作用} & \textbf{关键设计} \\
\hline
Devbox & 开发环境 & AWS EC2 预装源码和服务，预热池分配，启动约 10 秒，“牲口不是宠物” \\
\hline
编排状态机 & 流程控制 & 混合确定性节点，比如 lint、push，和 Agent 节点，比如实现功能、修 CI；该确定的地方确定，该灵活的地方灵活 \\
\hline
Toolshed MCP & 工具服务 & 集中式 MCP 服务，近 500 个工具，每个 Minion 拿到筛选后的子集 \\
\hline
反馈回路 & 质量保障 & Pre-push hook 秒级修 lint；推送后最多 2 轮 CI，覆盖 300 万+ 测试 \\
\hline
\end{tabular}
}



Stripe 的编排思路很像混合流水线。跑 lint、推送代码这类步骤走确定性流程；实现功能、修 CI 错误这类需要判断的部分交给 Agent。该死板的地方死板，该灵活的地方灵活。



他们还有一个理念：What's good for humans is good for agents。过去为人类工程师投入的 Devbox、工具链和开发者体验，在 Agent 上也会直接产生回报。Agent 不一定需要一套完全独立的基础设施，它更应该被当作开发环境中的一等公民。



Minions 底层是 Block 开源项目 \href{https://github.com/block/goose}{goose} 的一个 fork，Stripe 针对无人值守场景做了定制。



\subsection*{Mitchell Hashimoto：一个人的 Harness 工程学}


Mitchell Hashimoto 是 Vagrant、Terraform、Ghostty 终端模拟器的作者。他的路线和 Stripe 很不一样。他坚持一次只跑一个 Agent，并且保持深度参与。他明确说过：“我不打算跑多个 Agent，也不想跑。”



他的实践可以拆成六步：


{\footnotesize
\begin{tabular}{|p{\dimexpr 0.035\columnwidth-2\tabcolsep\relax}|p{\dimexpr 0.211\columnwidth-2\tabcolsep\relax}|p{\dimexpr 0.754\columnwidth-2\tabcolsep\relax}|}
\hline
\textbf{步骤} & \textbf{名称} & \textbf{做法} \\
\hline
1 & 放弃聊天模式 & 让 Agent 在能读文件、跑程序、发 HTTP 请求的环境里直接干活 \\
\hline
2 & 复现自己的工作 & 每件事做两次，一次自己做，一次让 Agent 做，他形容这个过程“痛苦至极” \\
\hline
3 & 下班前启动 Agent & 每天最后 30 分钟给 Agent 布置任务，比如深度调研、模糊探索、Issue 分拣 \\
\hline
4 & 外包确定性任务 & 挑出 Agent 几乎一定能做好的任务后台跑，建议关掉桌面通知，避免上下文切换 \\
\hline
5 & 工程化 Harness & Agent 每犯一次错，就工程化一个方案，尽量让它以后不再犯同类错误 \\
\hline
6 & 始终有 Agent 在跑 & 目标是 10-20\\% 的工作时间有后台 Agent 运行 \\
\hline
\end{tabular}
}



Ghostty 项目里的 \texttt{AGENTS.md} 很有代表性。每一行都对应一个过去的 Agent 失败案例。它是一个持续积累的防错系统。Agent 犯了一个新类型错误，就加一条规则，后面同类问题就能少一些。



\begin{figure}[H]
\centering
\includegraphics[width=0.85\textwidth]{/home/zhang/demos/latex_demos/Agent知识/HTML/一文搞懂 Harness Engineering：六层架构、上下文管理与一线团队实战 _ JavaGuide_files/continuously-evolving-harness-error-prevention-feedback-loop.pdf}
\caption*{持续进化的 Harness 防错反馈闭环}
\end{figure}



\subsection*{Birgitta Böckeler 对 Harness 的梳理}


Birgitta Böckeler 是 Thoughtworks 的 Distinguished Engineer，她在 Martin Fowler 网站上对 OpenAI 实践做过结构化分析。她更关心这些做法可以归到哪几类，以及还有哪些空白。



她把 Harness 组件归为三类：


{\footnotesize
\begin{tabular}{|p{\dimexpr 0.266\columnwidth-2\tabcolsep\relax}|p{\dimexpr 0.213\columnwidth-2\tabcolsep\relax}|p{\dimexpr 0.521\columnwidth-2\tabcolsep\relax}|}
\hline
\textbf{归类} & \textbf{关注点} & \textbf{典型实践} \\
\hline
Context Engineering & 管理 Agent 看到什么、什么时候看到 & 从巨大 \href{http://agents.md/}{AGENTS.md} 演化为入口文件 + 分层文档 \\
\hline
Architectural Constraints & 确保 Agent 不跑偏 & 自定义 Linter、结构测试、LLM Agent 充当约束 \\
\hline
Garbage Collection & 对抗熵积累 & 定期运行清理 Agent，扫描不一致和违规 \\
\hline
\end{tabular}
}



Böckeler 还提了几个判断，我觉得比案例本身更值得关注。



她认为 Harness 可能会变成新的服务模板。很多组织只有两三个主要技术栈，未来团队可能会从一组预制 Harness 中选择，就像今天从服务模板里创建新服务一样。



棕地项目改造会是最大挑战。公开成功案例大多是绿地项目，而把一个十年历史、没有清晰架构约束的代码库接入 Harness，要难得多。她把它比作在从没用过静态分析工具的代码库上运行静态分析，结果很可能是被警报淹没。她还提出 Ambient Affordances 这个概念：环境本身的结构特性会影响 Harness 能做多好。比如强类型语言天然有类型检查作为 sensor，清晰模块边界方便定义架构约束，Spring 这类框架也会抽象掉很多细节。



还有一个容易被忽略的问题：功能验证体系还很薄。现在很多讨论都集中在架构约束和熵管理上，但功能正确性验证仍然不够。Böckeler 的观察比较尖锐：很多团队让 AI 生成测试，再用这些测试验证 AI 生成的代码。这样做仍然缺少独立验证视角，她的原话是 puts a lot of faith into AI-generated tests, that's not good enough yet。



把这些案例放在一起看，共性比差异更明显：上下文污染、代码熵积累、工具调用可靠性，这三道坎几乎都会遇到。团队规模是 3 人还是 300 人，问题不太一样，但底层风险差不多。区别在于，有的团队等 Agent 出问题后再补救，有的团队一开始就把约束、验证和清理机制放进 Harness 里。后者的补救成本通常低很多。


\end{document}
